\documentclass{situationist-preprint}

\title{Libet's Experiment and the Decision to Participate}
\msid{ACT-05-b}
\msfamily{ACT-05}
\msdiscipline{Action, Agency and Free Will}
\mstype{Companion manuscript}
\mskeywords{Libet; readiness potential; decision to participate; picking; reasons-responsiveness; neuroscience of free will}

\begin{document}
\maketitle

\begin{abstract}
ACT-05 argues that the readiness potential precedes the decision to flex but not the decision to take part in the experiment, and that the latter is where the reasons bearing on the flexing operate. We develop the second claim. The decision to take part is a decision to adopt a policy of arbitrary picking, and each flexing executes that policy. We show that the reasons-responsiveness of the flexing is located in the policy by a counterfactual test, and we answer two objections: that the decision to take part has its own unconscious precursors, and that experiments on deliberate choice remove the distinction.
\end{abstract}

\section{The policy and its execution}

In Libet's protocol, participants agree to flex their wrists at a moment of their choosing within each trial and to report, by a clock, when they first became aware of the decision or urge to flex \citep{Libet1983}. ACT-05-a distinguishes the decision to take part (Definition~2.1) from the decision to flex (Definition~2.2), defines the locus of an action's reasons-responsiveness as the decision that would have been different had the reasons bearing on the action been different (Definition~2.3), and argues that the locus of the flexing is the decision to take part (Assumption~2.6).

The decision to take part has a distinctive content. It is not a decision to flex at a particular moment, but to flex once per trial at some moment, with the choice of moment left to chance or whim. In the terms of \citet{UllmannMargalitMorgenbesser1977}, it is a decision to pick rather than choose. Each flexing executes this policy, and the decision to flex now is the policy's way of resolving a question that the agent has deliberately left without reasons.

\begin{proposition}\label{p:test}
For each reason bearing on the flexing, the flexing would have varied with it through the decision to take part, and not through the decision to flex.
\end{proposition}

\begin{proof}
The reasons bearing on the flexing are of two kinds. Reasons for flexing at all, such as the wish to help the experimenters or the payment, would, had they been absent or outweighed, have led the participant not to take part, and so not to flex. Reasons bearing on when to flex within a trial are, by the design of the protocol, absent, since the instructions make every moment equally good. So any variation in the reasons that would have made a difference to the flexing makes it through the decision to take part. By Definition~2.3 of ACT-05-a, that decision is the locus.
\end{proof}

The participant's mechanism is responsive to reasons, and Libet's data concern the timing of the moments at which it has none.

\section{Objections and replies}

\begin{objection}[Precursors all the way down]
The decision to take part was itself preceded by neural activity of which the participant was not aware. If precursors threaten the decision to flex, they threaten it too.
\end{objection}

\begin{reply}
ACT-05-a grants Principle~P2, that an unconscious precursor preceding the locus would make the mechanism non-responsive, only to show that Libet's argument fails even on that premise. Extended to every decision, the principle would have no connection with the experiment: it would follow from the platitude that decisions have neural realizers. What the readiness potential was offered as evidence for is the specific claim that an unconscious process, not the agent's reasons, settles when she acts \citep{Libet1985}. That claim is true of the timing of the flexing, which the protocol made reasonless, and is not shown to be true of any decision in which reasons operate.
\end{reply}

\begin{objection}[Deliberate decisions]
Later studies record precursors of decisions with consequences, which the participants take for reasons.
\end{objection}

\begin{reply}
The study that compares the two kinds of decision directly found a readiness potential before arbitrary choices and no comparable potential before deliberate ones \citep{Maoz2019}. The distinction ACT-05 draws, between picking and choosing, thus appears in the neural data themselves. Predictive patterns found several seconds before reported decisions \citep{Soon2008} were obtained for arbitrary choices between two buttons, and predicted them only modestly above chance; they fall under Proposition~\ref{p:test} as Libet's data do.
\end{reply}

\section{Conclusion}

The decision to take part in Libet's experiment is a decision to pick, and the reasons bearing on the flexing operate through it. The readiness potential precedes the executions of that policy, not the policy. This supports the conclusion of ACT-05: the readiness potential precedes the decision to flex, not the decision to take part in the experiment, and neuroscience of this kind does not refute free will.

\begin{thebibliography}{9}
\bibitem[Libet(1985)]{Libet1985} Libet, B. (1985). Unconscious cerebral initiative and the role of conscious will in voluntary action. \emph{Behavioral and Brain Sciences}, 8(4), 529--566.
\bibitem[Libet et~al.(1983)]{Libet1983} Libet, B., Gleason, C.~A., Wright, E.~W. and Pearl, D.~K. (1983). Time of conscious intention to act in relation to onset of cerebral activity (readiness-potential): the unconscious initiation of a freely voluntary act. \emph{Brain}, 106(3), 623--642.
\bibitem[Maoz et~al.(2019)]{Maoz2019} Maoz, U., Yaffe, G., Koch, C. and Mudrik, L. (2019). Neural precursors of decisions that matter: an ERP study of deliberate and arbitrary choice. \emph{eLife}, 8, e39787.
\bibitem[Soon et~al.(2008)]{Soon2008} Soon, C.~S., Brass, M., Heinze, H.-J. and Haynes, J.-D. (2008). Unconscious determinants of free decisions in the human brain. \emph{Nature Neuroscience}, 11(5), 543--545.
\bibitem[Ullmann-Margalit and Morgenbesser(1977)]{UllmannMargalitMorgenbesser1977} Ullmann-Margalit, E. and Morgenbesser, S. (1977). Picking and choosing. \emph{Social Research}, 44(4), 757--785.
\end{thebibliography}

\end{document}
